\subsection{The one idea that explains everything else}\label{th01_what_edm_is-the-one-idea-that-explains-everything-else}

A pop song is designed to be listened to. A dance record is designed to be \emph{played}, by a
DJ, into another dance record, in a room full of people who are moving.

That single difference generates almost every formal property of the genre. Once you see
it, the conventions stop looking arbitrary.

Because a DJ has to beatmatch two records, the tempo has to be constant. No rubato, no
ritardando, no tempo changes for drama. A house record locks to one number for six minutes.

Because a DJ has to blend the end of one record into the start of the next, tracks begin
and end with long stretches of nothing but drums. Those are not filler. They are the
handles. A record that starts with a full arrangement in bar 1 is hard to mix and DJs will
quietly stop playing it.

Because people are dancing, the low frequencies carry the physical information and the
rhythm has to be legible from across a loud room. That is why the kick drum is so
prominent, why it usually lands on every beat, and why the sound design leans on things you
can feel in your chest.

Because a dancer is not tracking a narrative, the song does not need a verse, a chorus and
a bridge. It needs a slow accumulation of energy, a release, and enough variation to keep
people from noticing that eight minutes have passed. Repetition is the medium, not a
failure of imagination. The interest lives in what changes underneath the repetition.

Because the tracks are played in sequence over hours, an individual record does not have to
be self-sufficient. It can be a tool. A lot of the best techno is boring on headphones and
devastating at 2am in a sequence, and that is not a flaw in the record.

\begin{intuitionbox}
\textbf{The point.} When you cannot decide whether something in your own track is a good idea,
ask what it does to a person who is dancing and to the DJ who has to mix out of it. Nearly
every convention in the genre falls out of those two readers.
\end{intuitionbox}

\subsection{The materials}\label{th01_what_edm_is-the-materials}

Almost all dance music is made from four kinds of thing.

\textbf{Drums.} Usually synthesised or sampled from drum machines rather than recorded from a
kit. The Roland TR-808, TR-909 and TR-707 are the reference sounds for most of the genre.
They are not imitations of acoustic drums; they are their own instruments that happen to
occupy the same role.

\textbf{Bass.} One monophonic line in the lowest register, doing the harmonic and rhythmic
heavy lifting together. In a lot of house the bass is more important than the chords.

\textbf{Chords, pads and stabs.} The harmonic layer. In house this is very often an electric
piano, an organ, a string patch or a short filtered chord "stab" that plays on an offbeat.

\textbf{Everything else.} Vocals, samples, percussion, noise sweeps, sound effects. The
decoration and the identity.

There is no orchestra and no arranger. One person builds all four layers, usually in the
same afternoon, usually on a computer.

\subsection{Repetition, and why it works}\label{th01_what_edm_is-repetition-and-why-it-works}

The thing that alienates most newcomers is that a house track will play the same two bars
forty times. Coming from rock or classical, this reads as poverty.

It is closer to how you would think about a texture than how you would think about a
melody. When a pattern repeats long enough, your attention stops tracking the pattern and
starts tracking the differences: a hat that opens up, a filter slowly brightening, a new
percussion element sliding in. Very small changes become enormous. Removing the kick for
four bars, which is nothing on paper, is one of the most dramatic events available in the
genre.

This is also why producers obsess over sound quality to an extent that seems neurotic. If
the same clap is going to hit 400 times, it had better be a clap worth hearing 400 times.

Steve Reich and Terry Riley got here first, from a completely different direction, and if
you ever want to feel better about the intellectual respectability of loops, that lineage
is real.

\subsection{The two clocks}\label{th01_what_edm_is-the-two-clocks}

Dance music runs on two timescales at once and you need to hear both.

The fast clock is the bar. Four beats, about two seconds at house tempo. This is where
groove lives, and it is the level a dancer\textquotesingle s body is locked to.

The slow clock is the phrase. Eight bars, sixteen bars, thirty-two bars. Around fifteen,
thirty and sixty seconds respectively at house tempo. This is where the arrangement lives.
Changes happen on phrase boundaries almost without exception, and when you learn to feel
the eight-bar cycle you will start predicting every change in every track you hear. It is a
slightly annoying superpower.

Producers who cannot feel the slow clock make tracks that feel like they are wandering.
Almost every "my arrangement sounds wrong" problem is a phrase problem.

\subsection{What "EDM" means, and the vocabulary trap}\label{th01_what_edm_is-what-edm-means-and-the-vocabulary-trap}

The abbreviation is genuinely ambiguous and you should know which sense someone means.

In its broad sense, EDM is all of it. House, techno, garage, jungle, drum and bass,
dubstep, trance, everything with a programmed beat aimed at a dance floor. Used this way it
is just a category label, roughly equivalent to "rock".

In its narrow, American, post-2010 sense, "EDM" means the festival main-stage sound
specifically: big-room house, the Beatport top ten circa 2013, Ultra and EDC, the drop as
the entire point of the record. Avicii, Skrillex, Martin Garrix.

Most people who make house in Chicago use the second sense, and they use it as a mild
insult. If you tell a Smartbar regular you are getting into EDM they will hear something
different from what you meant. Say "house", or "dance music", or name the genre. This is
not gatekeeping so much as the ordinary situation where an outsider term and an insider
term coexist.

For the record: the music this folder is about is the ancestor of the festival stuff, not a
rival to it. House came out of Chicago in 1984 and went to New York, London, Ibiza, Berlin
and eventually back to America as a stadium product. Same river, much further downstream.

\subsection{What makes a genre a genre}\label{th01_what_edm_is-what-makes-a-genre-a-genre}

Genres in dance music are distinguished by a small number of parameters, and once you know
which parameters to check you can identify almost anything.

Tempo is the first thing, and it is a surprisingly hard constraint. House is 118 to 128
beats per minute. Techno is 125 to 145. Drum and bass is 170 to 175, and there is nothing
much in between 145 and 165 because those tempos are physically awkward to dance to.

The kick pattern is the second. Four to the floor (a kick on every beat) means house,
techno, trance, disco. A syncopated kick with a heavy backbeat means breakbeat, jungle,
garage, footwork.

Then swing, the sound palette, the harmonic density, the treatment of vocals, and the sound
system the music was designed for. \texttt{08\_genre\_map.md} lays all of this out in a table.

\subsection{Where to go next}\label{th01_what_edm_is-where-to-go-next}

\texttt{03\_rhythm.md} if you want to start making noise immediately. \texttt{02\_chicago\_house.md} if you
want to know whose shoulders you are standing on. Both are short.
